\section{Exact cleanup from a signed residual}
\label{sec:cleanup}

Define the note-scoped weighted absolute residual mass
\begin{equation}
 \mathcal M(r):=\sum_{i\in V}\sqrt{d_i}|r_i|.
 \label{eq:weighted-residual-mass}
\end{equation}
This is an analysis quantity, not a stopping rule.

\begin{lemma}[Signed-mass contraction of an exact push]
\label{lem:exact-mass}
For an arbitrary signed residual $r$, one exact $\omega=1$ push at $u$ gives
\begin{equation}
 \mathcal M(r^+)
 \leq\mathcal M(r)
 -\frac{2\alpha}{1+\alpha}\sqrt{d_u}|r_u|.
 \label{eq:exact-mass-drop}
\end{equation}
If $|r_u|/\sqrt{d_u}>\tau$ and the charged cost is $1+d_u$, then
\begin{equation}
 \mathcal M(r)-\mathcal M(r^+)
 >\frac{\alpha}{1+\alpha}\tau(1+d_u).
 \label{eq:exact-charge-drop}
\end{equation}
\end{lemma}

\begin{proof}
The exact push sets $r_u^+=0$.  For every neighbor $v$ it adds
\[
 \frac{1-\alpha}{1+\alpha}
 \frac{r_u}{\sqrt{d_ud_v}}
\]
to $r_v$.  The triangle inequality therefore bounds the total weighted
neighbor increase by
\[
 \sum_{v\in\mathcal N(u)}\sqrt{d_v}
 \frac{1-\alpha}{1+\alpha}
 \frac{|r_u|}{\sqrt{d_ud_v}}
 =\frac{1-\alpha}{1+\alpha}\sqrt{d_u}|r_u|.
\]
Subtract the removed self-mass.  Finally use
$\sqrt{d_u}|r_u|>\tau d_u$ and $1+d_u\leq2d_u$.
\end{proof}

\begin{corollary}[Work of a pure exact cleanup]
\label{cor:exact-cleanup-work}
If a sequential exact phase pushes only coordinates with
$|r_u|/\sqrt{d_u}>\tau$, then its charged work satisfies
\begin{equation}
 \Work_{\rm exact}
 \leq\frac{1+\alpha}{\alpha\tau}\mathcal M(r^{\rm start}).
 \label{eq:exact-cleanup-work}
\end{equation}
\end{corollary}

The bound proves that exact cleanup is robust to signed residuals, but its
explicit $1/\alpha$ also shows why cleanup alone is not an acceleration
proof.  The entering mass must be small, telescope across pairs, or be paid by
new frontier volume.
